% TOP_PROBLEM: {"problemId": "problem.uniqueness-of-tangent-cones-for-area-minimizing-integral-currents-in-higher-codimension", "canonicalTitle": "Uniqueness of Tangent Cones for Area-Minimizing Integral Currents in Higher Codimension", "releaseRank": 418, "primaryDomain": "analysis_pde", "primaryDomainLabel": "Analysis & PDE", "displayStatus": "open_with_solved_subcases", "releaseStatus": "open_with_solved_subcases"}
% !TeX program = lualatex
% Definition and sources only; this file is not an LLM solution attempt.
\documentclass[11pt]{article}
\usepackage[margin=25mm]{geometry}
\usepackage{fontspec}
\setmainfont{Latin Modern Roman}
\usepackage{amsmath,amssymb,mathtools}
\usepackage{unicode-math}
\setmathfont{Latin Modern Math}
\usepackage{xurl,hyperref}
\hypersetup{colorlinks=true,urlcolor=blue}
\setcounter{secnumdepth}{0}
\setlength{\parindent}{0pt}
\setlength{\parskip}{5pt}
\setlength{\emergencystretch}{3em}
\begin{document}
\section*{\texorpdfstring{Uniqueness of Tangent Cones for Area-Minimizing Integral Currents in Higher Codimension}{Uniqueness of Tangent Cones for Area-Minimizing Integral Currents in Higher Codimension}}\label{418}
\subsection{Definitions and mathematical statement}
An integral $m$-current represents an oriented rectifiable set with integer multiplicity and locally finite mass, with integral boundary. It is area-minimizing if no compactly supported competitor with the same boundary has smaller mass. For an interior point $p\in\operatorname{spt}T\setminus\operatorname{spt}\partial T$ and $r>0$, define
\[
 T_{p,r}=(\eta_{p,r})_\#T,
 \qquad\eta_{p,r}(x)=\frac{x-p}{r}.
\]
A tangent cone is a local weak limit along some $r_j\downarrow0$; weak convergence tests against compactly supported smooth differential forms. For $m\ge3$ and codimension $n\ge2$, the conjecture asserts that every such subsequential limit is the same current at every interior point. It is uniqueness of the full oriented multiplicity-bearing tangent current, not merely of its support or density, and concerns minimizing currents rather than arbitrary stationary varifolds.
\par\smallskip\textit{Formulation and concept references:} \href{https://proceedings.centre-mersenne.org/item/10.5802/slsedp.175.pdf}{[S1]}.

\subsection{Short English statement}
At every interior point of an area-minimizing integral current in higher codimension, do all infinitesimal blow-up sequences converge to the same tangent cone?
\subsection{Sources}
\begin{itemize}
\item[S1]Skorobogatova, Singularities of area-minimizing integral currents: recent progress and open directions, Seminaire Laurent Schwartz EDP et applications (2025).
\url{https://proceedings.centre-mersenne.org/item/10.5802/slsedp.175.pdf}.
\textit{Formulation source.}
\end{itemize}
\end{document}
